\documentclass[11pt]{article}
\usepackage[a4paper,margin=1in]{geometry}
\usepackage{amsmath,amssymb,amsthm}
\usepackage{booktabs}
\usepackage{array}
\usepackage{hyperref}

\newtheorem{definition}{Definition}
\newtheorem{corollary}{Corollary}
\newtheorem{proposition}{Proposition}

\newcommand{\XiBC}{\Xi_{\mathrm{BC}}}
\newcommand{\XiC}{\Xi_C}
\newcommand{\XiGamma}{\Xi_{\Gamma}}
\newcommand{\XiWD}{\Xi_{\mathrm{WD}}}
\newcommand{\CRE}{\mathrm{CRE}}
\newcommand{\CRCFT}{\mathrm{CRCFT}}
\newcommand{\CTMT}{\mathrm{CTMT}}

\title{Step 189: Framework Bridge Impossibility Corollary}
\author{RH Membrane Adequacy Track}
\date{}

\begin{document}
\maketitle

\section*{Verdict}

\[
\boxed{\mathrm{corollary\_verdict}=V_{\mathrm{corollary\_partial}}.}
\]

The target-entailment bridge impossibility corollary is proved from the Step
187 Dichotomy and the definition of \(\CRE\).  The stronger claim that every
such bridge is automatically \(\CRCFT\)-TE requires an exact/minimal bridge
hypothesis.  Without that hypothesis, the composite is forced into \(\CRCFT\),
but its mode can be TE, \(\CTMT\), or BF.

\section*{Definitions}

\begin{definition}[Typed bridge to the Riemann residual]
Let \(C\) be a non-\(\CRE\) RH-analogous carrier with established native
closure \(\XiC=0\).  A typed bridge \(B:C\to\XiBC\) is a legal collection of
carrier maps, response maps, energy comparisons, and defect records such that
closure of \(B\), together with \(\XiC=0\), derives
\[
\XiBC=0.
\]
\end{definition}

\begin{definition}[Exact/minimal bridge obligation]
Such a bridge is exact/minimal when the closure condition on \(B\) is
equivalent to closure of the enriched composite carrier \(C+B\).  Informally,
there is no additional independent obstruction outside the bridge record.
\end{definition}

\section*{Corollary}

\begin{corollary}[Framework bridge impossibility for non-CRE to Riemann]
Let \(C\) be a non-\(\CRE\) RH-analogous carrier with native closure
\(\XiC=0\).  Let \(B:C\to\XiBC\) be a typed bridge whose closure derives
\(\XiBC=0\) from \(\XiC=0\).  Then the enriched composite carrier \(C+B\) is
\(\CRE\), and hence falls under \(\CRCFT\) by the Step 187 Framework
RH-Carrier Dichotomy.

Consequently, a non-\(\CRE\) native closure cannot be transported to Riemann
RH as a free non-circular bridge under the framework.  The bridge is itself
the Riemann-strength obstruction.
\end{corollary}

\begin{proof}
The native carrier \(C\) is non-\(\CRE\), and by hypothesis its own residual
\(\XiC\) is already closed.  The bridge hypothesis says that if \(B\) closes,
then \(\XiC=0\) transports to
\[
\XiBC=0.
\]
The residual \(\XiBC\) is the active Riemann-zeta residual target in the
framework.  Therefore closure of the enriched composite \(C+B\) is a
Riemann-strength closure obligation.

By the Step 187 definition, any typed carrier whose closure is logically a
Riemann RH closure target is classically-RH-equivalent.  Hence \(C+B\) is
\(\CRE\).  The Step 187 Dichotomy then applies: every \(\CRE\) carrier falls
into \(\CRCFT\), in one of the modes \(\CTMT\), TE, or BF.

Since \(C\) was already closed natively, the newly introduced obstruction is
localized in the bridge data \(B\).  Thus no inherited non-\(\CRE\) closure can
serve as a free transport to \(\XiBC=0\); the transport itself is the
Riemann-strength task.
\end{proof}

\begin{proposition}[Exact-bridge TE form]
If, in addition, \(B\) is exact/minimal, then \(B\)-closure is
\(\CRCFT\)-TE: closing \(B\) is target-equivalent to closing the Riemann
residual through the composite \(C+B\).
\end{proposition}

\begin{proof}
By exactness, \(B\)-closure is equivalent to closure of the composite
\(C+B\).  By the preceding corollary, closure of \(C+B\) is the
Riemann-strength closure target.  Hence the bridge obligation is the
target-equivalent condition in the composite.  This is precisely the
\(\CRCFT\)-TE mode for an exact bridge.
\end{proof}

\paragraph{Why the unrestricted TE form is not forced.}
The Dichotomy only says that a \(\CRE\) composite lies in \(\CRCFT\).  It does
not say the mode is always TE.  If the bridge construction itself reduces to a
matrix-element family, the mode can be \(\CTMT\); if the bridge cannot compare
the ledgers, the mode can be BF.  Therefore the unrestricted statement
``every bridge is TE'' would require an extra exactness/minimality assumption.

\section*{Specializations}

\subsection*{Selberg/Maass to Riemann}

Step 185 established the native closure
\[
\XiGamma=0
\]
for the Selberg/Maass carrier by the full Selberg trace ledger.  A bridge
from this closure to \(\XiBC=0\) would enrich the Selberg carrier into a
Riemann-strength composite.  By the corollary, that composite is \(\CRE\) and
therefore \(\CRCFT\).  If the bridge obligation is exact/minimal, it is
\(\CRCFT\)-TE.

\subsection*{Weil--Deligne to Riemann}

Step 186 established the native closure
\[
\XiWD=0
\]
by Deligne purity and the Grothendieck--Lefschetz trace formula.  A typed
function-field-to-number-field bridge deriving \(\XiBC=0\) from this closure
would likewise make the enriched carrier \(\CRE\).  The same
\(\CRCFT\)-classification applies, with exact bridges landing in TE.

\section*{Strategic Implication for Riemann RH}

The Step 187 Dichotomy plus this bridge corollary leaves three framework
paths:
\begin{enumerate}
\item \textbf{Path 1.} Resolve a \(\CRCFT\)-stuck record on a \(\CRE\) carrier, such as the
Burnol projected-kernel \(\kappa\) theorem for Branch A/B, or a valid
replacement after the Branch C numerical foreclosure.
\item \textbf{Path 2.} Build a non-\(\CRE\)-to-Riemann bridge.  This is not free transport: by
the corollary it becomes a \(\CRE/\CRCFT\) obligation, and exact bridge closure
is TE.
\item \textbf{Path 3.} Find an in-scope carrier refuting the Dichotomy coverage conjecture.
No such instance was found through the RMT audit at Step 188.
\end{enumerate}

\section*{Nonclaim Boundary}

This corollary does not prove RH and does not close \(\XiBC\).  It does not
transfer Selberg or Weil--Deligne closure to Riemann RH.  It only shows that
any such transfer bridge becomes the Riemann-strength obstruction under the
framework.

\end{document}
